\section{Finite variation of the complete physical current}
\label{sec:v73-finite-current}
We prove that the original full current is a finite measure at each
collision count. The argument uses the complete finite-count height
bound, not a bound on separately differentiated word charts. It also
avoids an unproved definability assertion for a smooth cutoff.

\subsection{Thresholds on the unchanged finite collision graph}
The physical guard in \eqref{eq:v48-physical-decision-guards} is
\[
 H^\varepsilon=\prod_{j=0}^m\Theta\left(\frac{c_j}{\varepsilon e_j}\right)
       \prod_{j=0}^{m-1}\Theta\left(\frac{\Delta_j}{\varepsilon e'_j}\right),
 \quad e_j=q_0^{\min(j,m-j)/4},\quad
 e'_j=q_0^{\min(j,m-1-j)/4},\quad q_0=47/53.
\]
Here $\Theta$ is the original fixed smooth cutoff, equal to zero
on $(-\infty,1]$ and one on $[2,\infty)$. Its values lie in $[0,1]$.
Let $N=2m+1$ and enumerate the displayed normalized margins as
$a_1,\ldots,a_N$. Put $\sigma=d\Theta$ and
$V_\Theta=\|\sigma\|_{\TV}<\infty$. This is a finite signed measure
of mass one, supported in $[1,2]$. The proof does not require
monotonicity; for a nondecreasing cutoff $V_\Theta=1$.
For $y\in[1,2]^N$ define the positive threshold source
\begin{equation}\label{eq:v73-threshold-source}
 W_y(x)=1-\prod_{i=1}^N\one_{\{a_i(x)>y_i\}},\qquad
 b_y(t)\,dt=L_\#(W_y\mu).
\end{equation}

\begin{lemma}[Exact threshold representation]
\label{lem:v73-threshold-mixture}
The original complete source, with its original smooth guard, satisfies
\begin{equation}\label{eq:v73-threshold-mixture}
 1-H^\varepsilon=\int W_y\,d\sigma^{\otimes N}(y),\qquad
 b=\int b_y\,d\sigma^{\otimes N}(y)\quad\hbox{in }L^1(\R).
\end{equation}
Each $b_y$ is nonnegative and is dominated by the full exact-label
density. Uniformly in $R,\varepsilon,y,n,k$ at this count,
\begin{equation}\label{eq:v73-threshold-height}
 0\le b_y\le L_m\quad\hbox{a.e.},\qquad L_m=C A^m/c_*.
\end{equation}
\end{lemma}
\begin{proof}
The fundamental theorem of calculus gives
$\Theta(a)=\int\one_{\{y<a\}}\,d\sigma(y)$ for every real $a$.
Multiplication and the mass-one identity for $\sigma^{\otimes N}$
give the first equality. Finite signed Fubini applies because
$0\le W_y\le1$ and $\mu$ is finite. It gives the equality of
pushforward measures. The common kernel of
Lemma~\ref{lem:v73-common-kernel}, or joint Radon--Nikodym versions,
gives the displayed $L^1$ identity. Finally
Theorem~\ref{thm:v44-complete-density-bound} bounds every positive
restriction of the complete source by \eqref{eq:v73-threshold-height}.
The thresholds are auxiliary integration parameters, not a random
perturbation or a change to an actual orbit.
\end{proof}

\begin{lemma}[A uniform finite number of monotonicity intervals]
\label{lem:v73-threshold-monotonicity}
For each fixed $m$ there is an integer $J_m<\infty$ such that every
$b_y$ has a representative which is monotone on at most $J_m$
open intervals and has at most $J_m$ intervening point cells.
The same integer works for $R\in I$, all positive $\varepsilon$,
all thresholds, and every attainable exact label at this count.
No growth bound on $J_m$ as $m\to\infty$ is asserted.
\end{lemma}
\begin{proof}
Use the rational collision charts and the auxiliary physical graph
of Lemma~\ref{lem:v44-level-complexity}. At fixed $m$ there are
finitely many physical words, charts and exact labels. Circle,
flight and reflection equations, positive incidence and flight
length, and exclusion of an earlier hit are semialgebraic conditions
on this graph. Clearance is the minimum over the finite physical
candidate list, with the branch inequalities retained. Adding
$a_i>y_i$ therefore gives a semialgebraic threshold condition.
The original section is a finite union of rectangular chart pieces;
its initial, occupation and terminal tests are finite Boolean
conditions on this same graph. Displacement and occupation are fixed
at $(k,n)$, not summed in the pushforward under consideration.

The integration variables are the two initial coordinates, not the
ambient variables of the auxiliary graph. In a rational angular
chart $\alpha=\alpha_*+2\arctan x$, the initial density is a
constant multiple of $(1+x^2)^{-1}dx\,dp$, with the section factor
$c_*^{-1}$ inserted once. The unique physical collision graph and
projection show that its integrand, restricted by the exact event,
the threshold source and $L\le t$, is globally subanalytic. The
finite candidate alphabet follows from the uniform finite horizon.
Positive denominators and first-hit comparisons are retained under
projection; no continuation through a singularity is used.

It follows from the parameterized integration theorem
\cite[Theorem 1.3]{CMInt} that
\[
 M(R,\varepsilon,y,t)=\int W_y\one_{\{L\le t\}}\,d\mu
\]
is constructible, after summing the finitely many disjoint initial
chart pieces. Constructible functions are definable in
$\R_{\mathrm{an},\exp}$. By \eqref{eq:v73-threshold-height}, every
$t$-section of $M$ is $L_m$-Lipschitz. Define $g=\partial_tM$ where
the finite derivative exists and $g=0$ otherwise. The definition of
a derivative by quantified inequalities shows that $g$ is definable
in that same o-minimal structure. We need no assertion that a
constructible derivative is itself constructible. Absolute
continuity gives $g=b_y$ almost everywhere, and $0\le g\le L_m$
everywhere with this convention.

Uniform cell decomposition and the one-variable monotonicity theorem
in an o-minimal structure \cite{Coste} give a finite partition, with
the roof coordinate last, into cells on which these sections are
continuous and monotone or constant. The number of cells is uniform
in the displayed parameters for the fixed formula at this $m$.
Taking the maximum over the finitely many labels and chart unions
gives $J_m$. This is a finiteness statement for each fixed formula,
not a bound uniform over formulas of increasing collision length.
\end{proof}

\begin{theorem}[The complete scalar current is a finite measure]
\label{thm:v73-full-source-bv}
For every fixed $m\ge2$, the original full density
$b=b^\varepsilon_{n,k,m,R}$ belongs to $BV(\R)$. Its canonical
current of Theorem~\ref{thm:v72-complete-current} is a finite signed
measure, and for some finite $J_m$ as above,
\begin{equation}\label{eq:v73-full-bv-bound}
 \|Db\|_{\TV}\le 4(J_m+1)L_m V_\Theta^{2m+1}=:C_m.
\end{equation}
The bound is uniform in $R$, $\varepsilon$ and exact labels at this
fixed count. The current is compactly supported and has total mass
zero. It is the derivative of the complete original source, including
first incidence and original clearance outside the selected disks.
\end{theorem}
\begin{proof}
Each monotone interval in Lemma~\ref{lem:v73-threshold-monotonicity}
has variation at most $L_m$, and its traces lie in $[0,L_m]$.
Joining the finitely many pieces and extending by zero outside the
finite physical roof range gives
$\|Db_y\|_{\TV}\le4(J_m+1)L_m$. Isolated choices of representative
do not affect the distributional derivative. For
$\phi\in C_c^\infty(\R)$, signed Fubini and
\eqref{eq:v73-threshold-mixture} yield
\[
 |\langle Db,\phi\rangle|
 \le\int |\langle Db_y,\phi\rangle|\,d|\sigma|^{\otimes N}(y)
 \le4(J_m+1)L_m V_\Theta^N\|\phi\|_\infty.
\]
The order-zero representation theorem identifies $Db$ with a finite
signed measure and proves \eqref{eq:v73-full-bv-bound}; equivalently
this is the distributional characterization of $BV$ \cite{AFP}.
It is the same canonical derivative as before because both act as
$-\int b\phi'$. A smooth function equal to one on the physical roof
range shows that this derivative has total mass zero.
\end{proof}

The theorem establishes a property of the full physical sum. It does
not assert that a source multiplied by an arbitrary bounded mark,
nor an allocated subsource in \eqref{eq:v73-common-allocation-weights},
has a finite derivative. Nor does \eqref{eq:v73-full-bv-bound}
control $m^2\|Db\|_{\TV}$ as $m$ grows. The height $L_m$ is
exponential and the monotonicity number $J_m$ has no growth estimate
here. Finite variation supplies legitimate Jordan parts for the
next argument; their collision-uniform concentration remains a
separate quantitative question.
